\section{Discussion}
\label{sec:discussion}

This study shows that the reliability of connectome-based autism classification cannot be judged from discrimination alone. In ABIDE-I, the primary ensemble maintains similar AUROC under pooled and site-held-out validation, while additional analyses reveal important effects of cohort composition, calibration, regularisation, and the scale of model interpretation. Training-only recalibration corrects the average probability offset without reducing discrimination, although site-level calibration remains dependent on local autistic prevalence. Model explanations are more reproducible at the functional-system level, including across site-disjoint splits, than at the level of individual connections. Lower somatomotor connectivity persists after motion matching, while default-mode systems account for the largest overall share of classifier separation. Cognitive composition also influences measured performance, whereas classifier scores do not provide a corrected measure of symptom severity within autistic participants. Finally, the evaluated linear configurations perform better than the fixed tree ensembles, and the apparent benefit of windowed connectivity disappears after covariance regularisation is matched. Together, these findings support interpreting connectome classifiers through discrimination, calibration, cohort characteristics, analytical design, and reproducibility rather than through performance alone.

\subsection{Interpretation of the results}
\label{subsec:discussion-results}

Discrimination is maintained when an entire acquisition site is withheld from training. The primary ensemble achieves participant-pooled AUROCs of 0.776 under pooled ten-fold validation and 0.774 under LOSO validation (Table~\ref{tab:performance}), indicating little loss in discrimination when the test site is excluded from model fitting. Site-level performance is more variable, with AUROCs ranging from 0.589 to 0.912. The descriptive random-effects mean is 0.803, but it gives unequal weight to sites and is based on overlapping training data, so it \reviewdel{should not be interpreted}\reviewhl{is not interpreted} as another estimate of the participant-pooled AUROC or as a validated prediction for a new centre. In the wider QC cohort, accuracy is 0.706, which is comparable with previously reported ABIDE-I values of 0.700--0.708 \citep{Heinsfeldetal2018,Eslamietal2019,AlmuqhimSaeed2021}. The present contribution therefore lies in comparable discrimination under stricter validation and reliability analysis rather than numerical superiority.

Measured performance also varies with cohort composition. In unmatched autistic IQ bands, AUROC decreases from 0.873 to 0.719. Within autistic participants, classifier score is negatively associated with full-scale IQ, while the corresponding association is small in non-autistic participants. After IQ matching, the upper quartiles no longer show a monotonic decline, indicating that part of the original gradient reflects between-group IQ differences. However, diagnosis-by-IQ interactions remain significant for full-scale and verbal IQ. The evidence therefore supports sensitivity to cognitive composition rather than a simple IQ-dependent detectability gradient.

Motion-controlled analyses similarly refine the biological interpretation. Lower somatomotor connectivity persists after nuisance adjustment and becomes stronger after within-site matching on recorded mean FD. Under Freedman--Lane permutation, reduced connectivity is observed in 27 of 36 system pairs in the matched sample compared with 14 of 36 in the full cohort. The mean regional somatomotor effect also becomes more negative after matching. These findings are consistent with earlier reports \citep{Chenetal2015,Holigaetal2019,Ilioskaetal2023}, but they remain exploratory and do not establish a biomarker.

Several analytical choices have a substantial effect on reported performance. Selecting 2,000 of 19,900 connections using all diagnostic labels before cross-validation raises AUROC from 0.756 to 0.987, illustrating the inflation produced by outcome-informed feature selection. In contrast, region-level selection has only a small effect. \addhl{This analysis is a methodological check rather than a study contribution, and the two comparisons differ in cohort, shrinkage, and ranking statistic, so they do not isolate feature granularity alone.} On identical static tangent features and folds, the evaluated linear configurations achieve AUROCs of 0.735--0.757, compared with 0.611--0.624 for fixed-configuration tree ensembles. Because only logistic regression is tuned, this supports the tested linear configurations rather than a general superiority of linear model families. The apparent benefit of windowed connectivity is also removed when covariance shrinkage is matched: the windowed representation differs from the matched static baseline by only $-0.002$ AUROC. This shows that comparisons between static and windowed representations \reviewdel{should control for}\reviewhl{require control of} regularisation before attributing gains to temporal information.

The symptom analyses clarify what the classifier score represents. Several score--symptom associations are present when both diagnostic groups are analysed together, but none survives correction within autistic participants. For SRS total, the correlation decreases from 0.438 across both groups to 0.050 within autistic participants. The score therefore primarily reflects diagnostic-group separation rather than symptom severity within autism.

Calibration provides a further distinction between discrimination and probability reliability. The native probabilities show a small average offset despite calibration slopes close to one. This offset is consistent with penalisation of the intercept as a major contributor and is removed by training-only intercept recalibration without materially changing AUROC. Recalibration does, however, alter the balance between sensitivity and specificity at a fixed threshold. Residual calibration differences remain across withheld sites, and these offsets closely follow the autistic proportion of the site. \addhl{They are no longer heterogeneous across sites} when that base rate is known accurately, indicating that calibration at a new site depends on the local \reviewdel{prevalence}\reviewhl{autistic proportion}. By contrast, thresholds chosen from the training data \reviewdel{retain their intended sensitivity or specificity on unseen sites}\reviewhl{meet their predefined target in participant-pooled held-out predictions under both validation protocols}. These operating-point results apply to the present cohort, in which 46\% of participants are autistic, and \reviewdel{should not be interpreted}\reviewhl{do not constitute} evidence for population screening.

Model interpretation is also more reliable at broader biological scales. Default-mode pairs account for the largest overall share of classifier separation, while the somatomotor system contributes most per connection. System-level structure coefficients \reviewdel{reproduce}\reviewhl{are reproducible} across non-overlapping participant halves and across halves that share no acquisition site, whereas individual connections are less stable. Removing the highest-ranked system pairs reduces held-out AUROC more than random removal, showing that these systems contribute meaningfully to the fitted model while information remains distributed across the connectome.

Overall, the findings show that reliable connectome-based autism classification depends not only on discrimination, but also on cohort composition, motion control, regularisation, calibration, and the scale at which model explanations are interpreted.

\subsection{Limitations}
\label{subsec:limitations}

Several factors define the scope and generalisability of these findings. The primary cohort reflects the demographic structure of ABIDE-I, spanning 6.5 to 58 years with a median age of 14.8 years and substantially more male than female participants. The comparison group is predominantly typically developing rather than clinically referred, so the reported discrimination reflects separation within a research cohort rather than differential diagnosis. Clinically referred comparison groups are not available in ABIDE-I; operating points are therefore evaluated at the cohort autistic proportion of 46\%, while predictive values at lower \reviewdel{prevalences}\reviewhl{autistic proportions} are modelled rather than directly observed. LOSO validation tests transfer across acquisition sites within ABIDE-I, but does not establish generalisation to independent cohorts, scanners, populations, or preprocessing pipelines. The functional-connectivity findings are also conditional on the selected preprocessing strategy, CC200 parcellation, nuisance model, and absence of global signal regression. Motion-matched and regional analyses are therefore interpreted as exploratory, and agreement between group-level effects and classifier patterns is descriptive rather than causal. Classifier comparisons are likewise limited to the evaluated configurations because logistic regression is tuned while the tree ensembles use fixed settings. Additional sensitivity analyses using alternative atlases and reduced region sets do not materially improve performance, but help define the conditions under which the main results are obtained. Calibration and attribution analyses are interpreted according to the evidential status assigned to them in the Methods, and site-level calibration estimates apply only to the sites observed here. Structure coefficients \reviewdel{may still reflect}\reviewhl{potentially reflect} participant or site characteristics not explicitly removed from the classifier, while connections involving four regions with intermittent constant signals require cautious interpretation. Reproducibility is assessed using non-overlapping participant halves and site-disjoint splits, which provide stricter internal tests than overlapping folds; however, the 40 site divisions reuse the same 20 ABIDE-I sites and therefore do not substitute for validation in an independent cohort. Overall, these considerations mainly limit external generalisability and biological interpretation rather than the internal consistency of the findings. The somatomotor and default-mode patterns are therefore presented as cohort-level associations that warrant independent confirmation rather than as established biomarkers.

\subsection{Implications for reliability and generalisability}
\label{subsec:trustworthiness}

The findings indicate that reliability in connectome-based autism classification depends on several distinct properties rather than on discrimination alone. Within ABIDE-I, training-only recalibration corrects the average probability offset while preserving discrimination, supporting the reliability of the probability scale at the cohort level. Transfer to a new site is more conditional. Residual calibration differences are strongly related to the autistic proportion of the site, and the analyses in Figures~\ref{fig:trust}a and~\ref{fig:trust}b show that these differences \reviewdel{can be substantially reduced when the site prevalence}\reviewhl{are substantially reduced when the site autistic proportion} is known with reasonable accuracy. In contrast, \reviewdel{the operating point is more stable across sites}\reviewhl{participant-pooled operating-point performance is more stable than site-level probability calibration}. Training-derived thresholds \reviewdel{retain their intended sensitivity or specificity on held-out sites}\reviewhl{meet their intended target in participant-pooled held-out predictions under both validation protocols}, and decision-curve analysis\addhl{, a sensitivity analysis with no defined clinical decision, shows a net benefit above the better default at threshold probabilities of 0.30 to 0.50 in this 46\% cohort} (Figures~\ref{fig:trust}c and~\ref{fig:trust}d). Interpretation also shows a clear scale dependence: functional-system patterns \reviewdel{reproduce across independent participant halves}\reviewhl{are reproducible across non-overlapping participant halves} and site-disjoint splits, whereas individual connections are less stable. At the same time, variation with cognitive composition shows that discrimination is not independent of the population being evaluated. Taken together, the ensemble satisfies several important reliability criteria within ABIDE-I, particularly for discrimination, average calibration, \reviewdel{operating-point transfer}\reviewhl{aggregate held-out operating-point performance}, and system-level interpretation. Its use at a new site \reviewdel{will nevertheless require}\reviewhl{requires} attention to the local \reviewdel{prevalence}\reviewhl{autistic proportion} and cohort composition, while clinical generalisability \reviewdel{will require}\reviewhl{requires} validation in an independent dataset and a clinically representative population.

\subsection{Additional analyses and remaining requirements}
\label{subsec:required-validation}

Additional analyses clarify the uncertainty estimates, nesting of the design, and cognitive-ability findings. Site-clustered bootstrap intervals are wider than the DeLong and Wilson intervals, particularly at their lower bounds (Section~\ref{subsec:results-performance}), indicating that the parametric intervals understate some of the uncertainty associated with site clustering. However, all AUROC intervals remain above 0.5. Calibration is also evaluated using several measures rather than ECE alone (Section~\ref{subsec:results-calibration}), because the binned error does not fully capture the average probability offset. The cognitive-ability result is tested directly through diagnosis-by-IQ interactions. These interactions are negative under both validation protocols and survive correction for full-scale and verbal IQ, but not for performance IQ (Section~\ref{subsec:results-cognitive-ability}). The interpretation is therefore limited to full-scale and verbal IQ, and all analysed cohort counts are reconciled with explicit inclusion rules from the released records. The main results are also insensitive to stricter nesting. Refitting the tangent reference and feature scaler within every inner split changes the selected penalty in only 2 of 30 member fits under pooled validation and 2 of 60 under LOSO validation, while participant-pooled AUROC changes by less than 0.001 (Section~\ref{subsec:results-performance}). Similarly, refitting the nuisance model using only outer-training participants changes the region-selection AUROC by 0.0035 on average (Section~\ref{subsec:results-selection-scope}). The group-difference tests use Freedman--Lane residual permutation, and the within-site permutation of residualised connections is reported as a sensitivity analysis that identifies a subset of the same pairs (Section~\ref{subsec:results-system-differences}). Independent external validation nevertheless remains necessary. The study is therefore interpreted as a retrospective analysis of an aggregated research cohort, without claims about clinical screening, deployment, or performance at a new acquisition site. Its screening-related implications are methodological and descriptive, concerning how classifier scores are calibrated, interpreted, and reported rather than supporting their clinical use.

\subsection{Future work}
\label{subsec:future-work}

Priorities for future work are broader validation, stronger generalisability, and more consistent reporting of connectome-based autism classification. Independent external validation is required on a cohort such as ABIDE-II, either by reproducing the present preprocessing pipeline or by explicitly evaluating transfer across preprocessing strategies. The influence of cognitive composition also requires examination in cohorts with more balanced IQ distributions, greater representation of female participants and younger children, and clinically referred comparison groups. The somatomotor and default-mode findings require reassessment under alternative denoising strategies, including global signal regression and more detailed motion modelling, while the contribution of regions affected by constant signals requires separate examination. Classifier comparisons require tuned nonlinear models under the same nested validation framework, allowing model complexity to be evaluated more fairly. Calibration and system-level attribution also require testing across different solvers, preprocessing pipelines, and independent datasets to determine how consistently these properties are retained, and the use of a known base rate for calibration at a new site requires testing prospectively. More broadly, future studies benefit from reports of discrimination together with calibration, operating thresholds, cohort characteristics, and the reproducibility of model explanations. This direction supports more transparent comparison across studies and helps establish which aspects of connectome-based autism classification remain stable across populations, sites, and analytical settings.
